\section{Resultados}

Primeiro, apresenta-se o comportamento temporal dos sensores, a avaliação
principal por vaso (LOSO) na base completa e as sensibilidades temporais.
Em seguida, são mostradas a busca
combinatória de entradas e a comparação entre classificadores na análise
exploratória com cinco coletas excluídas. Por fim, apresentam-se os
resultados do teste histórico separado por coleta, nos recortes de solo
seco e úmido, e a análise da importância das entradas. As métricas de
protocolos diferentes são relatadas separadamente.

\subsection{Análise das Respostas dos Sensores}

A Figura \ref{fig:resposta_temporal_sensores} apresenta
o comportamento temporal dos sinais dos sensores durante
uma coleta de um vaso inoculado com \textit{H. glycines}.
O gráfico destaca as etapas de aquecimento, obtenção
da amostra e desaturação dos sensores.
% Revisar o idioma dos rótulos da figura antes da submissão.

\begin{figure}[H]
    \centering
    \includegraphics[width=\textwidth]{figures/explicativa.png}
    \caption{Comportamento temporal dos sinais dos sensores
    durante uma coleta de um vaso inoculado com
    \textit{H. glycines}.}
    \label{fig:resposta_temporal_sensores}

    \vspace{0.1cm}
    {\footnotesize Fonte: Elaborado pelos autores (2026).}
\end{figure}

Durante a fase inicial, os sensores apresentaram
variações acentuadas nos sinais, seguidas de um
período de estabilização. Na etapa de obtenção
da amostra, observaram-se oscilações e diferenças
nas amplitudes das respostas dos sensores MQ.
Com o início da desaturação, os sinais de parte
dos sensores apresentaram redução, caracterizando
a transição para o encerramento do ciclo de coleta.

\subsection{Validação LOSO na base completa}
\label{subsec:resultados_entradas_fixas}

Na base completa, o classificador Extra Trees (ET) com as nove entradas
fixas identificou corretamente 9 de 10 vasos, 31 de 36 coletas e
55.136 de 63.492 leituras (86,84\%). O conjunto
de cinco MQ selecionado entre 63 subconjuntos, acompanhado das três
variáveis ambientais, acertou 55.916 leituras (88,07\%), 32 coletas e
9 vasos. O resultado com nove entradas constitui a referência sem
seleção de subconjunto de sensores. A diferença de 1,23 ponto percentual
entre essas duas configurações não isola o efeito da seleção, pois as
entradas também mudam. Nas 63 combinações MQ mais ambiente, a acurácia
por leitura variou de 82,47\% a 88,07\%, com mediana de 86,39\%.

Na unidade biológica, o ET com nove entradas identificou corretamente
9 de 10 vasos (90\%; intervalo binomial exato de 95\%: 55,50--99,75\%).
Foram 5 de 5 vasos do grupo inoculado (intervalo: 47,82--100\%) e
4 de 5 controles (intervalo: 28,36--99,49\%). O único erro por vaso
ocorreu em um controle. Esses intervalos são amplos porque há apenas
cinco vasos por grupo; as 63.492 leituras repetidas não ampliam o número
de unidades biológicas independentes.

No mesmo ET de nove entradas, a acurácia balanceada foi de 87,08\% por
leitura, 86,22\% por coleta e 90,00\% por vaso. Os intervalos descritivos
de 95\% obtidos por 10.000 reamostragens estratificadas de vasos foram,
respectivamente, 76,78--95,84\%, 75,08--97,06\% e 70,00--100\%.
Cada reamostragem preservou todas as previsões do vaso sorteado; os
intervalos são condicionais às entradas e ao modelo já fixados.

O ET com apenas as três variáveis ambientais acertou 53.787 leituras
(84,71\%), 29 coletas e 10 vasos; com apenas os seis MQ, 50.041 leituras
(78,81\%), 28 coletas e 9 vasos. A diferença média de acurácia por
leitura dentro de vaso entre as nove entradas e somente ambiente foi de
2,80 pontos percentuais. A enumeração exploratória das 1.024 inversões
de sinal das dez diferenças por vaso resultou em $p=0{,}348$; por acurácia
por coleta dentro de vaso, $p=0{,}50$. Essas comparações não demonstram
ganho estável dos MQ além das variáveis ambientais.

A Tabela~\ref{tab:controles_loso_auc} reúne as configurações fixas da
base completa no processamento originalmente consolidado e a regressão
logística. Esta última acertou 43.621 de
63.492 leituras, 24 de 36 coletas e 8 de 10 vasos. A AUC por coleta foi
calculada com as probabilidades médias de cada coleta e é apresentada
como medida descritiva das previsões LOSO reunidas.

\begin{table}[H]
    \centering
    \footnotesize
    \caption{Controles com entradas fixas na base completa e AUC por coleta.}
    \label{tab:controles_loso_auc}
    \begin{tabular}{p{0.18\textwidth}p{0.16\textwidth}p{0.25\textwidth}ccc}
        \hline
        Modelo & Entradas & Acertos/leituras (\%) & Coletas & Vasos & AUC por coleta \\
        \hline
        ET & Nove entradas & 55.136/63.492 (86,84) & 31/36 & 9/10 & 0,864 \\
        ET & Somente seis MQ & 50.041/63.492 (78,81) & 28/36 & 9/10 & 0,842 \\
        ET & Somente ambiente & 53.787/63.492 (84,71) & 29/36 & 10/10 & 0,873 \\
        Regressão logística & Nove entradas & 43.621/63.492 (68,70) & 24/36 & 8/10 & 0,752 \\
        \hline
    \end{tabular}

    \vspace{0.1cm}
    {\footnotesize Fonte: Elaborado pelos autores (2026).}
\end{table}

A regressão logística obteve acurácia balanceada de 68,37\% por leitura
e 66,56\% por coleta, com F1 macro de 68,41\% e 66,56\%, respectivamente.
As AUCs não estabelecem superioridade estatística
entre os modelos, pois combinam previsões de rodadas LOSO ajustadas em
conjuntos de treinamento diferentes.

Como verificação de sensibilidade, a validação LOSO foi repetida com o
índice de solo e a mediana de pressão calculados exclusivamente nos vasos
de treinamento de cada rodada. A Tabela~\ref{tab:preprocessamento_loso_treino}
compara a faixa estrita de pressão, definida pela mediana do treino
$\pm 0{,}50$~kPa, com a ausência dessa restrição adicional. Nos dois
cenários, manteve-se o filtro fixo de saltos de pressão.

\begin{table}[htbp]
    \centering
    \scriptsize
    \caption{Validação LOSO com pré-processamento ajustado no treino de cada rodada.}
    \label{tab:preprocessamento_loso_treino}
    \begin{tabular}{llrcc}
        \hline
        Faixa estrita & Entradas & Acertos/leituras (\%) & Coletas & Vasos \\
        \hline
        Sim & Seis MQ + ambiente & 55.136/63.492 (86,84) & 31/36 & 9/10 \\
        Sim & Somente ambiente & 53.788/63.492 (84,72) & 29/36 & 10/10 \\
        Sim & Somente seis MQ & 50.041/63.492 (78,81) & 28/36 & 9/10 \\
        Sim & Cinco MQ + ambiente\textsuperscript{a} & 55.916/63.492 (88,07) & 32/36 & 9/10 \\
        Não & Seis MQ + ambiente & 56.491/66.081 (85,49) & 32/36 & 9/10 \\
        Não & Somente ambiente & 56.036/66.081 (84,80) & 30/36 & 10/10 \\
        Não & Somente seis MQ & 51.670/66.081 (78,19) & 28/36 & 9/10 \\
        Não & Cinco MQ + ambiente\textsuperscript{a} & 57.972/66.081 (87,73) & 32/36 & 9/10 \\
        \hline
    \end{tabular}

    \vspace{0.1cm}
    {\footnotesize \textsuperscript{a}Subconjunto escolhido anteriormente
    nos mesmos vasos; resultado exploratório. Ambiente: índice de solo,
    temperatura e pressão.}
\end{table}

Com a faixa estrita ajustada no treino, o ET com as nove entradas manteve
55.136/63.492 acertos por leitura, 31/36 coletas e 9/10 vasos, os mesmos
totais da análise com parâmetros globais. No controle somente ambiente,
houve diferença de uma leitura correta: 53.788 na reanálise, contra
53.787 no processamento global. A mediana de pressão do treino
foi 93,69~kPa em nove rodadas e 93,68~kPa em uma. Os extremos de solo
mudaram em duas rodadas; sete leituras de teste ficaram fora da escala
$[0,1]$ obtida no respectivo treino. Sem a faixa estrita, permaneceram
2.589 leituras adicionais: a acurácia por leitura das nove entradas foi
85,49\%, com 32/36 coletas e 9/10 vasos corretos. Como os dois cenários
avaliam números diferentes de leituras, a diferença das acurácias por
leitura não constitui uma comparação pareada das mesmas amostras. A
diferença média de acurácia por leitura dentro de vaso entre as nove
entradas e somente ambiente foi de 2,80 pontos percentuais com a faixa
estrita ($p=0{,}348$) e de 1,56 ponto percentual sem ela ($p=0{,}660$),
em enumerações exploratórias de inversões de sinal por vaso. Esses dados
não demonstram ganho estável dos MQ além das variáveis ambientais. A
largura de $\pm 0{,}50$~kPa foi mantida do processamento originalmente
adotado e não foi validada em novos vasos; esta reanálise verifica a
sensibilidade dos resultados às duas regras avaliadas, sem constituir
validação externa do protocolo de filtragem.

\subsection{Sensibilidade ao dia, ao período e ao recorte temporal}
\label{subsec:resultados_sensibilidade_temporal}

O referencial que atribuiu a cada coleta a classe majoritária observada
no próprio dia acertou 24 de 36 coletas (66,67\%). Como utilizou os
rótulos avaliados, esse número é apenas descritivo. Quando a classe
majoritária do dia foi estimada somente nos outros nove vasos de cada
rodada LOSO, o acerto caiu para 14 de 36 coletas (38,89\%).

Na análise \textit{leave-one-date-out}, que reservou uma data inteira
para teste a cada rodada, o ET com nove entradas acertou 10 de 36
coletas com a faixa estrita e 13 de 36 sem ela
(Tabela~\ref{tab:validacao_data_fora}). Em contraste, no LOSO por vaso
com faixa estrita, 31 de 36 coletas foram identificadas corretamente.
Por leitura, a análise por data acertou 19.105 de 63.492 registros
(30,09\%) com a faixa estrita e 24.613 de 66.081 (37,25\%) sem essa
faixa. Os denominadores diferem porque a restrição de pressão altera
as leituras elegíveis.

\begin{table}[htbp]
    \centering
    \footnotesize
    \caption{Acertos por coleta do ET com nove entradas ao reservar cada data para teste.}
    \label{tab:validacao_data_fora}
    \begin{tabular}{ccc}
        \hline
        Dia de maio & Com faixa estrita & Sem faixa estrita \\
        \hline
        11 & 2/7 & 2/7 \\
        12 & 1/6 & 4/6 \\
        13 & 0/4 & 0/4 \\
        18 & 5/9 & 5/9 \\
        19 & 2/10 & 2/10 \\
        \hline
        Total & 10/36 & 13/36 \\
        \hline
    \end{tabular}
\end{table}

O dia 13 continha apenas coletas controle, e nenhuma das quatro foi
classificada corretamente quando essa data ficou fora do treinamento.
As datas anteriores e posteriores à irrigação também correspondem a
faixas distintas do índice de solo. Assim, a queda no teste por data é
compatível com dependência de condições temporais, ambientais ou
operacionais compartilhadas pelas coletas, sem permitir identificar
qual desses fatores a causou. Essa análise não avalia vasos inéditos,
pois os mesmos vasos podem aparecer em outras datas do treinamento.

A Tabela~\ref{tab:sensibilidade_periodos} apresenta a avaliação com
entradas fixas restrita aos dias 18 e 19 e as transferências entre os
períodos anteriores e posteriores à irrigação. Na transferência entre
períodos, vasos medidos em ambos puderam aparecer no treino e no teste.

\begin{table}[H]
    \centering
    \scriptsize
    \caption{Sensibilidade do ET ao período de aquisição, sem exclusão de coletas.}
    \label{tab:sensibilidade_periodos}
    \begin{tabular}{p{0.18\textwidth}p{0.15\textwidth}p{0.25\textwidth}p{0.08\textwidth}p{0.16\textwidth}}
        \hline
        Treino e teste & Entradas & Acertos/leituras (\%) & Coletas & Bal. por coleta (\%) \\
        \hline
        LOSO em 18--19 & Nove entradas & 22.249/27.365 (81,30) & 15/19 & 78,89 \\
                         & Seis MQ & 17.537/27.365 (64,09) & 13/19 & 68,33 \\
                         & Somente ambiente & 20.539/27.365 (75,06) & 14/19 & 73,89 \\
        \hline
        11--13 para 18--19 & Nove entradas & 13.459/27.365 (49,18) & 9/19 & 50,00 \\
                            & Seis MQ & 12.350/27.365 (45,13) & 8/19 & 43,33 \\
                            & Somente ambiente & 13.459/27.365 (49,18) & 9/19 & 50,00 \\
        \hline
        18--19 para 11--13 & Nove entradas & 13.117/36.127 (36,31) & 6/17 & 30,00 \\
                            & Seis MQ & 20.304/36.127 (56,20) & 10/17 & 60,71 \\
                            & Somente ambiente & 22.895/36.127 (63,37) & 11/17 & 57,14 \\
        \hline
    \end{tabular}

    \vspace{0.1cm}
    {\footnotesize Bal. = acurácia balanceada. Fonte: Elaborado pelos autores (2026).}
\end{table}

Com nove entradas, o acerto por coleta foi de 15/19 no LOSO restrito
aos dias 18 e 19, de 9/19 ao treinar nos dias 11 a 13 e testar nos dias
18 e 19, e de 6/17 no sentido inverso. A queda nas transferências é
compatível com mudança de condições entre períodos, mas o desenho não
isola o efeito da irrigação. A padronização dos seis MQ por dia, com
parâmetros estimados apenas nos vasos de treinamento, produziu
49.927/63.492 acertos por leitura (78,64\%) e 28/36 por coleta no LOSO
completo. Sem essa padronização, os seis MQ acertaram 50.041 leituras
(78,81\%) e as mesmas 28 coletas.

No recorte fixo de 200 a 800~s, o ET foi reajustado nas 36 coletas.
Permaneceram 51.888 leituras: 23.217 do grupo inoculado e 28.671 dos
controles. Com nove entradas, o modelo acertou 46.053/51.888 leituras
(88,75\%) e 32/36 coletas; com seis MQ, 41.252/51.888 (79,50\%) e
27/36; e com somente ambiente, 44.312/51.888 (85,40\%) e 31/36.
A duração efetivamente retida por coleta teve mediana de 309,00~s
(181,88--548,32~s) no grupo inoculado e 345,69~s
(260,32--587,33~s) nos controles. Como a abertura da válvula não foi
marcada individualmente, esse intervalo é um recorte temporal de
sensibilidade e não identifica a fase de exposição ao \textit{headspace}
em todas as coletas.

\subsection{Busca combinatória na base após exclusões}
\label{subsec:resultados_busca_combinatoria_filtrada}

Na análise de sensibilidade, foram excluídas cinco coletas, totalizando
7.546 leituras; restaram 31 coletas e 55.946 leituras dos dez vasos.
Essas coletas foram selecionadas após a inspeção dos erros e dos sinais,
sem comprovação de que todas as leituras excluídas eram defeituosas.
Apenas uma coleta excluída pertencia à classe minoritária do próprio dia:
o controle do vaso 2 em 12 de maio. As outras quatro ocorreram em 19 de
maio, quando havia cinco coletas de cada grupo.

O ET avaliou os 63 subconjuntos não vazios dos seis MQ, isoladamente
e acompanhados de solo, temperatura e pressão. Para apresentar as
configurações sem MQ7, a Tabela~\ref{tab:busca_mq_filtrada} mostra o
primeiro colocado entre os 31 subconjuntos formados pelos outros cinco
MQ em cada cenário. Essa restrição de apresentação não altera os
resultados da busca original.
\begin{table}[H]
    \centering
    \footnotesize
    \caption{Melhores subconjuntos sem MQ7 no LOSO exploratório com 31 coletas.}
    \label{tab:busca_mq_filtrada}
    \begin{tabular}{p{0.13\textwidth}p{0.43\textwidth}p{0.21\textwidth}p{0.12\textwidth}}
        \hline
        Cenário & Entradas do primeiro colocado & Acertos/leituras & Acurácia \\
        \hline
        Só MQ & MQ2, MQ3, MQ8, MQ135 e MQ138 & 53.015/55.946 & 94,76\% \\
        MQ + ambiente & MQ2, MQ3, MQ8, MQ135 e MQ138 + solo, temperatura e pressão & 54.783/55.946 & 97,92\% \\
        \hline
    \end{tabular}
\end{table}

A busca mais ampla das 511 combinações individuais das nove variáveis
também colocou os cinco MQ mais solo, temperatura e pressão em primeiro
lugar, com 54.783 acertos em 55.946 leituras (97,92\%). O ranking
recalculado por acurácia balanceada manteve essa configuração em primeiro
lugar, com 97,88\%. A escolha e a avaliação dos subconjuntos usaram os
mesmos vasos após as exclusões; ambas as posições no ranking são
exploratórias e não estimam o desempenho em vasos novos.

\subsection{Comparação dos classificadores na base após exclusões}
\label{subsec:resultados_modelos_loso_filtrado}

Para comparar ET, Random Forest e TCN nas mesmas entradas, foram fixados
MQ2, MQ3, MQ8, MQ135 e MQ138, isoladamente e acompanhados de solo,
temperatura e pressão. A Tabela~\ref{tab:modelos_cinco_mq_filtrado}
apresenta os acertos por leitura após reunir as dez rodadas LOSO.
\begin{table}[H]
    \centering
    \caption{Desempenho por leitura dos três modelos com os mesmos cinco MQ, na base de 31 coletas.}
    \label{tab:modelos_cinco_mq_filtrado}
    \begin{tabular}{lcc}
        \hline
        Modelo & Cinco MQ & Cinco MQ + ambiente \\
        \hline
        Extra Trees & 53.015/55.946 (94,76\%) & 54.783/55.946 (97,92\%) \\
        Random Forest & 47.172/55.946 (84,32\%) & 53.385/55.946 (95,42\%) \\
        TCN & 49.352/55.946 (88,21\%) & 49.789/55.946 (88,99\%) \\
        \hline
    \end{tabular}
\end{table}
Os resultados desta comparação usam os mesmos cinco MQ apresentados na
Tabela~\ref{tab:busca_mq_filtrada}. Como as exclusões e a escolha das entradas foram
exploratórias, esses valores não estimam desempenho externo. Os três
modelos também tiveram estratégias de ponderação e arquiteturas
diferentes, de modo que a tabela não estabelece superioridade geral.

Na mesma base de 31 coletas, as AUCs por coleta para cinco MQ foram
1,000 no ET, 0,992 no RF e 0,971 na TCN. Com cinco MQ e as três
variáveis ambientais, foram 1,000, 1,000 e 0,983, respectivamente.
Essas AUCs foram reunidas das rodadas LOSO depois da seleção das entradas
e da exclusão exploratória das coletas. Portanto, não representam uma
avaliação independente das decisões de pré-processamento e seleção.

\subsection{Análise exploratória no recorte de solo seco}

As próximas subseções apresentam recortes exploratórios da partição
histórica por coleta, definidos por faixas de um índice operacional do solo.
Os modelos foram treinados uma vez e as previsões do mesmo teste foram
separadas posteriormente pelo índice. Quantidade e modo de irrigação e tempo
de estabilização não foram controlados para esta comparação. Além disso, os
mesmos vasos puderam aparecer em treinamento e teste em dias distintos;
essas métricas não devem ser comparadas diretamente às da validação por vaso.
% Conferir se os rótulos nas figuras históricas ainda dizem X-TCN; o texto
% usa TCN para coincidir com a arquitetura descrita em Materiais e Métodos.

No solo seco, a TCN apresentou o melhor desempenho nas três métricas, como
mostra a Tabela~\ref{tab:resultados_solo_seco_vacuo}.

\begin{table}[htbp]
    \centering
    \caption{Desempenho dos classificadores no solo seco com coleta a vácuo.}
    \label{tab:resultados_solo_seco_vacuo}
    \begin{tabular}{lrrr}
        \hline
        \textbf{Modelo} & \textbf{Acurácia (\%)} &
        \textbf{Acurácia balanceada (\%)} & \textbf{F1 macro (\%)} \\
        \hline
        a) TCN         & \textbf{98,47} & \textbf{98,03} & \textbf{98,38} \\
        b) Extra Trees & 96,15          & 95,06          & 95,87          \\
        \hline
    \end{tabular}

    \vspace{0.1cm}
    {\footnotesize Fonte: Elaborado pelos autores (2026).}
\end{table}

Em relação ao Extra Trees, a TCN apresentou ganhos de 2,32 pontos
percentuais em acurácia, 2,97 pontos percentuais em acurácia balanceada e 2,51
pontos percentuais em F1 macro. A rede classificou corretamente todas as
leituras do grupo controle e alcançou recall de 96,07\% nas leituras do
grupo inoculado. O ET também reconheceu todas as leituras do grupo controle,
mas obteve recall de 90,12\% para o grupo inoculado.

A Figura~\ref{fig:matrizes_confusao_solo_seco_vacuo} apresenta as matrizes de
confusão dos dois classificadores no solo seco.

\begin{figure}[H]
    \centering
    \includegraphics[width=\linewidth]{figures/matrizes_solo_seco.png}
    \caption{Matrizes de confusão no teste em solo seco: (a) TCN e
    (b) Extra Trees. As linhas representam as classes reais e as colunas,
    as classes previstas; os percentuais são relativos à classe real.}
    \label{fig:matrizes_confusao_solo_seco_vacuo}

    \vspace{0.1cm}
    {\footnotesize Fonte: Elaborado pelos autores (2026).}
\end{figure}

A TCN reduziu o número de erros de 643 para 256 leituras, correspondendo a
uma redução de 60,19\%. Nesse teste por coleta, a TCN apresentou menos
erros que o ET no recorte seco.

\subsection{Análise exploratória no recorte de solo úmido}

No recorte denominado solo úmido pelo índice operacional, o Extra Trees
teve métricas superiores às da TCN nas três medidas avaliadas,
conforme a Tabela~\ref{tab:resultados_solo_umido_vacuo}.

\begin{table}[htbp]
    \centering
    \caption{Desempenho dos classificadores no solo úmido com coleta a vácuo.}
    \label{tab:resultados_solo_umido_vacuo}
    \begin{tabular}{lrrr}
        \hline
        \textbf{Modelo} & \textbf{Acurácia (\%)} &
        \textbf{Acurácia balanceada (\%)} & \textbf{F1 macro (\%)} \\
        \hline
        Extra Trees & \textbf{74,46} & \textbf{62,96} & \textbf{63,33} \\
        TCN         & 63,58          & 41,52          & 38,87          \\
        \hline
    \end{tabular}

    \vspace{0.1cm}
    {\footnotesize Fonte: Elaborado pelos autores (2026).}
\end{table}

O ET obteve recall de 84,61\% para as leituras do grupo inoculado e de
41,31\% para as leituras do grupo controle. A diferença entre a acurácia de 74,46\%
e a acurácia balanceada de 62,96\% mostra que o resultado foi influenciado pelo
desbalanceamento entre as classes.

A TCN alcançou recall de 83,04\% para o grupo inoculado, mas não
reconheceu nenhuma leitura do grupo controle. Por esse motivo, apesar da
acurácia de 63,58\%, a acurácia balanceada caiu para 41,52\% e o F1 macro para
38,87\%. O resultado mostra que o padrão temporal identificado no solo seco
não se manteve estável no solo úmido.

A Figura~\ref{fig:matrizes_confusao_solo_umido_vacuo} apresenta as matrizes de
confusão dos dois classificadores no solo úmido.

\begin{figure}[H]
    \centering
    \includegraphics[width=\linewidth]{figures/matrizes_solo_umido.png}
    \caption{Matrizes de confusão no teste em solo úmido: (a) TCN e
    (b) Extra Trees. }
    \label{fig:matrizes_confusao_solo_umido_vacuo}

    \vspace{0.1cm}
    {\footnotesize Fonte: Elaborado pelos autores (2026).}
\end{figure}

\subsection{Comparação exploratória das faixas do índice de solo}

Os dois modelos apresentaram métricas superiores no recorte denominado solo
seco. As diferenças descritivas
entre as condições são apresentadas na Tabela
\ref{tab:comparacao_umidade_vacuo}.

\begin{table}[htbp]
    \centering
    \caption{Comparação do desempenho entre solo seco e solo úmido nas coletas a vácuo.}
    \label{tab:comparacao_umidade_vacuo}
    \footnotesize
    \begin{tabular}{llrrr}
        \hline
        \textbf{Modelo} & \textbf{Condição} & \textbf{Acurácia (\%)} &
        \textbf{Acurácia bal. (\%)} & \textbf{F1 macro (\%)} \\
        \hline
        Extra Trees & Solo seco  & 96,15 & 95,06 & 95,87 \\
                      & Solo úmido & 74,46 & 62,96 & 63,33 \\
                      & Diferença  & 21,69 & 32,10 & 32,54 \\
        \hline
        TCN         & Solo seco  & 98,47 & 98,03 & 98,38 \\
                      & Solo úmido & 63,58 & 41,52 & 38,87 \\
                      & Diferença  & 34,89 & 56,51 & 59,51 \\
        \hline
    \end{tabular}

    \vspace{0.1cm}
    {\footnotesize Fonte: Elaborado pelos autores (2026).}
\end{table}

No ET, as métricas do recorte úmido foram inferiores às do seco em 21,69
pontos percentuais de acurácia, 32,10 de acurácia balanceada e 32,54 de
F1 macro. Na TCN, as diferenças foram de 34,89, 56,51 e 59,51 pontos
percentuais, respectivamente. A queda das métricas balanceadas coincidiu
com a mudança da faixa operacional do índice de solo. Como também houve
mudança de data e irrigação entre os períodos, não é possível atribuir
essa queda isoladamente à umidade.

\subsection{Subamostragem balanceada da partição histórica}
\label{subsec:resultados_subamostragem_historica}

As previsões do ET no teste histórico foram reavaliadas após igualar o
número de leituras dos dois grupos em cada faixa do índice de solo,
sem reajustar o modelo. A Tabela~\ref{tab:subamostragem_historica}
apresenta o sorteio com semente 42 e, separadamente, as leituras
excedentes que não entraram no subconjunto balanceado.

\begin{table}[H]
    \centering
    \scriptsize
    \caption{Sensibilidade do ET à subamostragem balanceada no teste histórico.}
    \label{tab:subamostragem_historica}
    \begin{tabular}{p{0.16\textwidth}p{0.20\textwidth}p{0.20\textwidth}ccc}
        \hline
        Conjunto & Leituras por grupo: inoc./controle & Acertos/leituras & Acurácia (\%) & Bal. (\%) & F1 macro (\%) \\
        \hline
        Seco balanceado & 6.507/6.507 & 12.371/13.014 & 95,06 & 95,06 & 95,05 \\
        Úmido balanceado & 1.271/1.271 & 1.584/2.542 & 62,31 & 62,31 & 60,57 \\
        Ambos os recortes & 7.778/7.778 & 13.955/15.556 & 89,71 & 89,71 & 89,71 \\
        Leituras excedentes & 2.881/3.678 & 6.132/6.559 & 93,49 & 92,59 & 93,25 \\
        \hline
    \end{tabular}

    \vspace{0.1cm}
    {\footnotesize Bal. = acurácia balanceada. Fonte: Elaborado pelos autores (2026).}
\end{table}

No recorte seco balanceado, o ET manteve 95,06\% de acurácia
balanceada. No recorte úmido, obteve 62,31\%. O conjunto de leituras
excedentes reuniu controles do recorte seco e vasos inoculados do
recorte úmido; sua métrica não constitui teste independente, pois as
leituras provêm das mesmas coletas. Os valores balanceados dependem de
um único sorteio e não substituem os resultados do teste completo.

\subsection{Comparação dos modelos no teste separado por coleta}

A posição relativa dos modelos variou entre os recortes. No solo seco, a
TCN apresentou os melhores resultados, alcançando 98,47\% de acurácia,
98,03\% de acurácia balanceada e 98,38\% de F1 macro. No solo úmido, o ET
apresentou as maiores métricas, com 74,46\%, 62,96\% e 63,33\%, respectivamente.

Ao considerar conjuntamente todas as leituras do teste histórico, o ET
atingiu 90,83\% de acurácia e a TCN atingiu 89,91\%. Essa diferença
global é pequena, mas oculta comportamentos distintos: a TCN foi mais eficaz
no solo seco e falhou na identificação do grupo controle no solo úmido,
enquanto o Extra Trees apresentou maior estabilidade entre as condições.

Esses resultados descrevem o desempenho dos modelos na partição
histórica por coleta. Como os recortes diferem também em data e condição
operacional, não permitem atribuir a mudança relativa de desempenho à
umidade isoladamente.

\subsection{Importância dos sensores e das variáveis ambientais no teste por coleta}


Para comparar a contribuição das entradas nos dois classificadores, foi
realizada uma análise adicional com os mesmos conjuntos de treinamento e teste,
separados por coleta. Em cada classificador foram ajustadas duas versões: uma
com os seis sinais MQ brutos (MQ2, MQ3, MQ7, MQ8, MQ135 e MQ138) e outra com
esses sinais acrescidos do índice de umidade do solo, da temperatura e da
pressão interna. Os sinais MQ brutos foram mantidos nas duas versões porque os
MQ corrigidos utilizados na análise anterior já incorporam uma compensação
estatística baseada nas variáveis ambientais. Assim, os resultados desta
ablação não substituem os percentuais apresentados nas tabelas anteriores.
O modelo arbóreo utilizou 900 árvores, seleção de atributos pela raiz
quadrada do número de entradas e mínimo de dez amostras por folha. A TCN
foi treinada por 60 épocas com janelas de 128 leituras e passo de 64; a
semente aleatória foi 42 nos dois casos.

O teste incluiu 22.115 leituras de 12 coletas, das quais 16.692 leituras de
oito coletas foram classificadas operacionalmente como solo seco e 5.423
leituras de quatro coletas como solo úmido. Cada modelo foi treinado uma única
vez com todas as coletas de treinamento; a separação por condição do solo foi
aplicada apenas à avaliação. Na TCN, a probabilidade atribuída a cada
leitura corresponde à média das probabilidades das janelas temporais que a
contêm. A Tabela~\ref{tab:ablacao_ambiente_vacuo} apresenta os resultados
obtidos por leitura.

\begin{table}[htbp]
    \centering
    \caption{Desempenho dos classificadores com e sem variáveis ambientais nas coletas a vácuo.}
    \label{tab:ablacao_ambiente_vacuo}
    \scriptsize
    \begin{tabular}{llrrrr}
        \hline
        \textbf{Modelo} & \textbf{Entradas} &
        \textbf{Acurácia total (\%)} &
        \textbf{Bal. total (\%)} &
        \textbf{Bal. seco (\%)} &
        \textbf{Bal. úmido (\%)} \\
        \hline
        Extra Trees & MQ              & 89,86 & 89,90 & 99,93 & 38,56 \\
        Extra Trees & MQ + ambiente   & 87,93 & 87,89 & 99,92 & 33,30 \\
        TCN         & MQ              & 80,06 & 79,77 & 91,50 & 26,07 \\
        TCN         & MQ + ambiente   & 89,41 & 89,54 & 95,12 & 46,15 \\
        \hline
    \end{tabular}

    \vspace{0.1cm}
    {\footnotesize Bal. = acurácia balanceada. Fonte: Elaborado pelos autores (2026).}
\end{table}

Para medir a importância das entradas nos dois modelos, foi adotada a
\emph{importância por permutação} (\emph{permutation feature importance}),
conforme o princípio de perturbar uma variável por vez e medir a alteração do
desempenho do modelo, conforme o procedimento de \cite{fisher2019modelreliance}.
Em cada repetição, os
valores de uma coluna foram embaralhados entre as leituras do teste completo,
mantendo-se o modelo treinado e as demais colunas. A importância corresponde à
diferença entre a acurácia balanceada original e a obtida após a permutação,
expressa em pontos percentuais. Foram realizadas três permutações por variável
em cada cenário. O mesmo procedimento e a mesma métrica foram usados para o
Extra Trees e para a TCN; nesta última, as janelas foram reconstruídas a
partir das leituras perturbadas antes da predição. A
Tabela~\ref{tab:importancia_variaveis_vacuo} apresenta a média das três quedas.

\begin{table}[htbp]
    \centering
    \caption{Importância por permutação no teste completo: queda da acurácia balanceada, em pontos percentuais.}
    \label{tab:importancia_variaveis_vacuo}
    \scriptsize
    \begin{tabular}{lrrrr}
        \hline
        \textbf{Variável} &
        \textbf{ET: MQ} &
        \textbf{TCN: MQ} &
        \textbf{ET: MQ + amb.} &
        \textbf{TCN: MQ + amb.} \\
        \hline
        Pressão interna & --    & --    & 7,13  & 23,49 \\
        Índice de solo & --    & --    & -0,37 & 8,12  \\
        Temperatura     & --    & --    & -2,60 & -1,01 \\
        MQ2             & 2,40  & -3,68 & -0,24 & 17,64 \\
        MQ3             & 3,84  & -8,92 & 1,13  & -3,77 \\
        MQ7             & 3,88  & -1,14 & 1,63  & 3,07  \\
        MQ8             & 9,28  & -5,11 & 1,96  & 5,25  \\
        MQ135           & 9,40  & 8,20  & 0,57  & 3,77  \\
        MQ138           & 0,17  & -5,13 & 1,43  & 10,59 \\
        \hline
    \end{tabular}

    \vspace{0.1cm}
    {\footnotesize -- = variável ausente do modelo. Fonte: Elaborado pelos autores (2026).}
\end{table}

Com as variáveis ambientais, a pressão interna produziu a maior queda de
acurácia balanceada nos dois modelos: 7,13 pontos percentuais no Extra Trees
e 23,49 na TCN. Na TCN, também se destacaram MQ2 (17,64), MQ138 (10,59)
e o índice de umidade do solo (8,12 pontos percentuais). Sem as entradas
ambientais, MQ135 apresentou a maior queda nos dois classificadores. Essas
diferenças mostram que a posição de um sensor no ranking depende das demais
entradas oferecidas ao modelo. Valores negativos indicam que, nessa avaliação,
o embaralhamento da variável aumentou a acurácia balanceada; não representam
uma influência física negativa do sensor sobre o grupo experimental.

A inclusão das variáveis ambientais elevou a acurácia balanceada total da
TCN de 79,77\% para 89,54\%, mas reduziu a do Extra Trees de 89,90\% para
87,89\%. No recorte úmido, a TCN passou de 26,07\% para 46,15\% e o ET
passou de 38,56\% para 33,30\%. Portanto, o ganho da TCN com as
informações ambientais não eliminou a dificuldade de classificação nessa
condição. A distinção entre solo seco e úmido foi feita por um índice
operacional do sensor de solo, sem medição calibrada da água em diferentes
profundidades. Como o teste úmido contém apenas quatro coletas e a condição de
umidade também varia com o dia de coleta, esses resultados não isolam um efeito
causal da água sobre os gases ou sobre os nematoides. A importância preditiva
da pressão também pode refletir características da câmara e da operação do
sistema de vácuo.

